\documentclass[11pt,a4paper]{article}

% ─── Tiếng Việt qua XeLaTeX ──────────────────────────────────────────────────
\usepackage{fontspec}
\usepackage{polyglossia}
\setmainlanguage{vietnamese}
\setmainfont{Latin Modern Roman}

% ─── Layout & tiện ích ───────────────────────────────────────────────────────
\usepackage[margin=2.5cm]{geometry}
\usepackage{booktabs}
\usepackage{xcolor}
\usepackage{listings}
\usepackage{enumitem}
\usepackage{fancyhdr}

% ─── Cấu hình listing cho code ───────────────────────────────────────────────
\definecolor{codebg}{rgb}{0.96,0.96,0.96}
\definecolor{codekw}{rgb}{0.13,0.29,0.53}
\definecolor{codecmt}{rgb}{0.30,0.50,0.30}
\definecolor{codestr}{rgb}{0.64,0.20,0.20}
\definecolor{framecol}{rgb}{0.70,0.70,0.70}

\lstset{
  basicstyle=\ttfamily\small,
  numbers=left,
  numberstyle=\tiny\color{gray},
  stepnumber=1,
  numbersep=8pt,
  backgroundcolor=\color{codebg},
  frame=single,
  rulecolor=\color{framecol},
  breaklines=true,
  breakatwhitespace=false,
  showstringspaces=false,
  keywordstyle=\color{codekw}\bfseries,
  commentstyle=\color{codecmt}\itshape,
  stringstyle=\color{codestr},
  tabsize=2,
  captionpos=b,
  literate=%
    {á}{{\'a}}1 {à}{{\`a}}1 {ả}{{\h a}}1 {ã}{{\~a}}1 {ạ}{{\d a}}1
    {â}{{\^a}}1 {ấ}{{\'{\^a}}}1 {ầ}{{\`{\^a}}}1
    {ă}{{\u a}}1 {ắ}{{\'{\u a}}}1 {ằ}{{\`{\u a}}}1
    {é}{{\'e}}1 {è}{{\`e}}1 {ế}{{\'{\^e}}}1 {ề}{{\`{\^e}}}1 {ệ}{{\d{\^e}}}1
    {í}{{\'i}}1 {ì}{{\`i}}1 {ó}{{\'o}}1 {ò}{{\`o}}1 {ô}{{\^o}}1 {ố}{{\'{\^o}}}1
    {ơ}{{\horn o}}1 {ớ}{{\'{\horn o}}}1 {ợ}{{\d{\horn o}}}1
    {ú}{{\'u}}1 {ù}{{\`u}}1 {ư}{{\horn u}}1 {ứ}{{\'{\horn u}}}1 {ự}{{\d{\horn u}}}1
    {đ}{{\dj}}1 {Đ}{{\DJ}}1 {ý}{{\'y}}1 {ị}{{\d i}}1 {ộ}{{\d{\^o}}}1
}

\lstdefinestyle{java}{language=Java}
\lstdefinestyle{sql}{language=SQL}

% ─── Header/Footer ───────────────────────────────────────────────────────────
\pagestyle{fancy}
\fancyhf{}
\lhead{\small TIP -- Money Tracker}
\rhead{\small Kiến trúc \& Thiết kế Hệ thống}
\cfoot{\thepage}
\renewcommand{\headrulewidth}{0.4pt}

% ─── hyperref (load gần cuối) ────────────────────────────────────────────────
\usepackage[colorlinks=true, linkcolor=blue, urlcolor=blue, citecolor=blue]{hyperref}

% ─── Metadata ────────────────────────────────────────────────────────────────
\title{\textbf{Kiến trúc \& Thiết kế Hệ thống}\\[4pt]
       \large TIP -- Tap Into Prosperity (Money Tracker)\\
       \normalsize Hỗ trợ đa thiết bị, đồng bộ Delta \& toàn vẹn số dư}
\author{Nhóm phát triển TIP}
\date{\today}

\begin{document}
\maketitle
\begin{abstract}
\noindent Báo cáo trình bày các quyết định kiến trúc cốt lõi của ứng dụng quản lý
tài chính cá nhân \textbf{TIP}: mô hình ba tầng \emph{offline-first}, cơ chế đồng bộ
\emph{Delta Sync} dựa trên con trỏ, xác thực \emph{JWT (HMAC-SHA)} kết hợp quản lý
phiên đa thiết bị, bảo đảm toàn vẹn số dư theo nguyên tắc \emph{số dư $=\Sigma$ giao dịch},
cô lập dữ liệu đa tài khoản và kiểm thử tích hợp end-to-end. Mọi mô tả bám sát mã nguồn
thực tế của dự án (package gốc \texttt{vn.edu.usth.tip}).
\end{abstract}

\tableofcontents
\newpage

% ═════════════════════════════════════════════════════════════════════════════
\section{Tổng quan kiến trúc ba tầng}
\label{sec:overview}

Hệ thống được tổ chức thành ba tầng rõ ràng, theo triết lý \textbf{Offline-First}:
cơ sở dữ liệu cục bộ \texttt{Room/SQLite} là nguồn sự thật cho trải nghiệm người dùng,
còn \texttt{PostgreSQL} trên cloud (Neon) là nguồn sự thật cho toàn hệ thống giữa nhiều
thiết bị.

\begin{enumerate}[leftmargin=2em]
  \item \textbf{Tầng Android (offline-first).} Kiến trúc \texttt{MVVM + Repository}:
        \texttt{Fragment/Activity} chỉ \texttt{observe} \texttt{LiveData} do \texttt{AppViewModel}
        (Financial Engine) phát ra; \texttt{Repository} điều phối giữa Room (cache cục bộ)
        và Retrofit (REST). Mọi thao tác CRUD ghi vào Room \emph{trước}, đồng bộ lên server
        \emph{sau} qua \texttt{WorkManager}.
  \item \textbf{Tầng Network/JWT.} \texttt{Retrofit/OkHttp} gắn header
        \texttt{Authorization: Bearer <token>} cho mỗi request, tự động làm mới token qua
        \texttt{Authenticator} và buộc đăng xuất khi gặp 401/403.
  \item \textbf{Tầng Spring Boot $\rightarrow$ PostgreSQL.} Kiến trúc ba lớp chuẩn
        \texttt{Controller} $\rightarrow$ \texttt{Service} $\rightarrow$ \texttt{JPA Repository},
        mọi thao tác đổi số dư được bọc trong \texttt{@Transactional} (ACID), ràng buộc
        \texttt{CHECK(amount > 0)} và khóa ngoại ở tầng database.
\end{enumerate}

\begin{center}
\fbox{\parbox{0.92\textwidth}{\centering\ttfamily\small
ANDROID (Java) \\
UI (Fragment/ViewModel) $\rightarrow$ Repository (offline-first) \\
Room SQLite (UX truth) \quad | \quad Retrofit/OkHttp (HTTP + JWT) \\[4pt]
$\downarrow$ HTTPS + JWT \\[4pt]
SPRING BOOT 4 / Java 21 \\
Controller $\rightarrow$ Service $\rightarrow$ JPA Repository \quad @Transactional \\[4pt]
$\downarrow$ JDBC \\[4pt]
NEON POSTGRESQL (system truth) \quad CHECK(amount>0), FK
}}
\end{center}

% ═════════════════════════════════════════════════════════════════════════════
\section{Đồng bộ: Delta Sync dựa trên con trỏ}
\label{sec:sync}

Thay vì tải lại toàn bộ dữ liệu mỗi lần đồng bộ, hệ thống chỉ truyền những bản ghi
\emph{đã thay đổi kể từ lần đồng bộ trước} (delta). Mỗi entity
(\texttt{transactions}, \texttt{accounts}, \texttt{categories}, \texttt{budgets},
\texttt{goals}, \texttt{debts}) có một endpoint \texttt{/delta} và một con trỏ (cursor)
riêng do \emph{server} cấp, lưu trong \texttt{SyncPrefs}.

\subsection{Ba tham số đảm bảo ``không sót, không trùng''}

\begin{itemize}[leftmargin=2em]
  \item \texttt{updatedSince} (cursor cũ): chỉ lấy bản ghi \texttt{updatedAt > since};
        lần đầu dùng \texttt{SyncConstants.EPOCH\_CURSOR}.
  \item \texttt{untilTimestamp} (snapshot): chốt \texttt{until = now()} ngay đầu phiên;
        mọi trang đều thấy cùng một ``ảnh chụp'' \texttt{updatedAt <= until}. Chính
        \texttt{until} trở thành con trỏ mới cho phiên kế tiếp.
  \item \texttt{lastUpdatedAt + lastId} (keyset pagination): phân trang theo cặp khóa
        \texttt{(updatedAt, id)} thay cho \texttt{OFFSET}; \texttt{id} phá thế hòa khi
        nhiều bản trùng \texttt{updatedAt} tới từng mili-giây.
\end{itemize}

Truy vấn delta dùng \texttt{Slice<T>} (\emph{không} \texttt{COUNT(*)}) -- client chỉ cần
biết ``còn nữa hay không'' (\texttt{hasMore = slice.hasNext()}):

\begin{lstlisting}[style=java,caption={Truy vấn delta keyset (TransactionRepository.findDelta)}]
@Query("SELECT t FROM Transaction t WHERE t.user.id = :uid " +
       "AND t.updatedAt > :since AND t.updatedAt <= :until " +     // cửa sổ (since, until]
       "AND (:lastTs IS NULL OR t.updatedAt > :lastTs " +          // keyset pagination
       "     OR (t.updatedAt = :lastTs AND t.id > :lastId)) " +
       "ORDER BY t.updatedAt ASC, t.id ASC")
Slice<Transaction> findDelta(uid, since, until, lastTs, lastId, pageable);
\end{lstlisting}

\subsection{Last-Write-Wins (LWW)}

Khi cùng một bản ghi bị sửa ở nhiều nơi, hệ thống dùng chiến lược \textbf{Last-Write-Wins}
dựa trên \texttt{updatedAt}: bản nào mới hơn sẽ thắng. Logic được xử lý null-safe để bản
ghi cũ chưa từng được ``stamp'' không vô tình ghi đè dữ liệu tốt. Đặc biệt, server trả về
\emph{cả} bản ghi đã soft-delete (\texttt{deletedAt != null}) để thiết bị khác biết mà xóa
theo (cơ chế \emph{tombstone}); nếu lọc bỏ, lệnh xóa sẽ không bao giờ lan truyền.

% ═════════════════════════════════════════════════════════════════════════════
\section{Xác thực JWT \& Quản lý phiên đa thiết bị}
\label{sec:auth}

Server chạy \textbf{stateless}: mỗi request tự mang danh tính qua một JWT đã ký bằng
\texttt{HMAC-SHA} với khóa \texttt{jwt.secret}. Để hỗ trợ đa thiết bị và thu hồi tức thì,
hệ thống bổ sung bảng \texttt{sessions} và claim \texttt{sid} (id phiên) trong access token.

\subsection{Vòng đời token}

\begin{itemize}[leftmargin=2em]
  \item \textbf{Đăng nhập/đăng ký} cấp \emph{access token} (JWT mang claim \texttt{sid},
        hạn \texttt{jwt.expiration} = 24h) + \emph{refresh token} opaque (lưu SHA-256,
        hạn \texttt{jwt.refresh-expiration} = 30 ngày).
  \item \textbf{Refresh xoay vòng (rotation).} Mỗi lần đổi refresh token, server sinh
        token mới và lưu lại SHA-256 mới $\rightarrow$ token cũ trở nên vô hiệu, chống
        replay.
  \item \textbf{Thu hồi tức thì.} \texttt{JwtAuthFilter} kiểm tra phiên \texttt{sid} còn
        sống ở \emph{mỗi} request; đăng xuất từ xa có hiệu lực ngay ở request kế tiếp.
\end{itemize}

\subsection{Access token mang claim \texttt{sid}}

\begin{lstlisting}[style=java,caption={JwtUtil -- phát access token mang claim ``sid'' (id phiên)}]
public String generateToken(UserDetails userDetails, UUID sessionId) {
    Map<String, Object> claims = new HashMap<>();
    if (sessionId != null) {
        claims.put("sid", sessionId.toString());   // gắn id phiên vào token
    }
    return createToken(claims, userDetails.getUsername());
}
\end{lstlisting}

\texttt{JwtAuthFilter} đọc \texttt{sid} và hỏi \texttt{SessionService.isActive(sid)}; nếu phiên
đã bị thu hồi/hết hạn thì \emph{không} thiết lập authentication, khiến Spring Security trả
401 ngay. Token cũ không mang \texttt{sid} vẫn được chấp nhận (tương thích ngược).

\subsection{Refresh token xoay vòng (SHA-256)}

Refresh token là chuỗi ngẫu nhiên 256-bit; server chỉ lưu \emph{băm} SHA-256 của nó, không
bao giờ lưu token thô. Mỗi lần \texttt{rotate()} sẽ kiểm tra hợp lệ rồi thay băm mới:

\begin{lstlisting}[style=java,caption={SessionService.rotate -- kiểm tra \& xoay refresh token}]
public SessionCreation rotate(String rawRefresh) {
    Session s = sessionRepository.findByRefreshTokenHash(sha256(rawRefresh))
            .orElseThrow(() -> new InvalidTokenException("Refresh token không hợp lệ"));
    if (s.isRevoked() || s.getExpiresAt().isBefore(OffsetDateTime.now())) {
        throw new InvalidTokenException("Phiên đã hết hạn hoặc bị thu hồi");
    }
    String newRaw = generateRawToken();        // XOAY: sinh token mới
    s.setRefreshTokenHash(sha256(newRaw));      // lưu băm mới -> token cũ vô hiệu
    s.setLastSeenAt(OffsetDateTime.now());
    s.setExpiresAt(OffsetDateTime.now().plus(refreshExpirationMs, ChronoUnit.MILLIS));
    sessionRepository.save(s);
    return new SessionCreation(s, newRaw);
}
\end{lstlisting}

\subsection{Tự động làm mới phía client (OkHttp Authenticator)}

Khi gặp 401, OkHttp gọi \texttt{Authenticator}: client dùng refresh token (qua client
\emph{không} xác thực để tránh đệ quy) lấy access token mới rồi thử lại request, giữ người
dùng đăng nhập mà không phải nhập lại. Một khóa (\texttt{refreshLock}) đảm bảo một loạt 401
đồng thời chỉ kích hoạt tối đa một lần refresh.

\begin{lstlisting}[style=java,caption={RetrofitClient -- OkHttp Authenticator tự refresh token}]
private static okhttp3.Authenticator tokenAuthenticator(TokenManager tokenManager) {
    return (route, response) -> {
        if (tokenManager == null) return null;
        if (responseCount(response) >= 2) return null;   // đã thử 1 lần -> ngừng
        String refresh = tokenManager.getRefreshToken();
        if (refresh == null) return null;
        synchronized (refreshLock) {
            String current = tokenManager.getToken();
            String tried   = stripBearer(response.request().header("Authorization"));
            if (current != null && !current.equals(tried)) {   // thread khác đã refresh
                return response.request().newBuilder()
                        .header("Authorization", "Bearer " + current).build();
            }
            retrofit2.Response<AuthResponse> r =
                    getAuthApi().refresh(new RefreshRequest(refresh)).execute();
            if (r.isSuccessful() && r.body() != null && r.body().getToken() != null) {
                tokenManager.updateTokens(r.body().getToken(), r.body().getRefreshToken());
                return response.request().newBuilder()
                        .header("Authorization", "Bearer " + r.body().getToken()).build();
            }
            return null;   // refresh hỏng -> 401 lan ra -> interceptor đăng xuất
        }
    };
}
\end{lstlisting}

\subsection{UI quản lý thiết bị}

Màn \texttt{SessionManagementFragment} (``Thiết bị đăng nhập'') liệt kê các phiên đang hoạt
động, đánh dấu thiết bị hiện tại, cho phép thu hồi một thiết bị hoặc ``đăng xuất các thiết bị
khác''. Các endpoint mới do \texttt{AuthController} và \texttt{AccountController} cung cấp được
tóm tắt ở Bảng~\ref{tab:endpoints}.

\begin{table}[ht]
\centering
\caption{Các endpoint mới phục vụ đa thiết bị \& toàn vẹn số dư}
\label{tab:endpoints}
\begin{tabular}{@{}llll@{}}
\toprule
\textbf{Method} & \textbf{Đường dẫn} & \textbf{Quyền} & \textbf{Chức năng} \\
\midrule
POST   & \texttt{/api/auth/login}                    & public & Đăng nhập, cấp access + refresh token \\
POST   & \texttt{/api/auth/register}                 & public & Đăng ký, tự tạo phiên đầu tiên \\
POST   & \texttt{/api/auth/refresh}                  & public & Xoay refresh token, cấp access mới \\
POST   & \texttt{/api/auth/logout}                   & auth   & Thu hồi phiên hiện tại \\
GET    & \texttt{/api/auth/sessions}                 & auth   & Liệt kê phiên đang hoạt động \\
DELETE & \texttt{/api/auth/sessions/\{id\}}          & auth   & Thu hồi một phiên (đăng xuất từ xa) \\
POST   & \texttt{/api/auth/sessions/logout-others}   & auth   & Thu hồi mọi phiên khác, giữ phiên này \\
\midrule
POST   & \texttt{/api/accounts/\{id\}/reconcile}     & auth   & Tính lại số dư một ví từ giao dịch \\
POST   & \texttt{/api/accounts/reconcile-all}        & auth   & Tính lại số dư tất cả ví \\
GET    & \texttt{/api/accounts/delta}                & auth   & Delta sync ví (cursor-based) \\
\bottomrule
\end{tabular}
\end{table}

\noindent Cấu hình \texttt{SecurityConfig}: chỉ \texttt{/api/auth/login|register|refresh} là
\texttt{permitAll}, mọi đường dẫn còn lại \texttt{.authenticated()}; \texttt{csrf} tắt,
\texttt{SessionCreationPolicy.STATELESS}, băm mật khẩu bằng BCrypt.

% ═════════════════════════════════════════════════════════════════════════════
\section{Toàn vẹn số dư: số dư $=\Sigma$ giao dịch}
\label{sec:balance}

Số dư là một bộ đếm cộng dồn nên \emph{không an toàn} dưới LWW khi nhiều thiết bị cùng sửa
(dễ ``trôi lệch''). Vì vậy hệ thống coi \textbf{số dư ví $=\Sigma(\text{thu}) - \Sigma(\text{chi}+\text{chuyển})$}
và tính lại sau \emph{mọi} thao tác create/update/delete và sau \emph{mỗi} batch \texttt{/sync},
thay vì cộng dồn.

\subsection{Tính lại số dư (recompute / reconcile)}

\begin{lstlisting}[style=java,caption={TransactionService.recomputeBalance \& truy vấn computeBalanceForAccount}]
// TransactionService -- số dư = Sigma giao dịch (authoritative, không trôi LWW)
private void recomputeBalance(Account account) {
    if (account == null) return;
    BigDecimal computed = transactionRepository.computeBalanceForAccount(account.getId());
    if (computed == null) computed = BigDecimal.ZERO;
    if (account.getBalance() == null || account.getBalance().compareTo(computed) != 0) {
        account.setBalance(computed);   // chỉ ghi khi lệch -> tránh bump updatedAt thừa
        accountRepository.save(account);
    }
}

// TransactionRepository -- tổng có dấu theo loại giao dịch, bỏ qua bản đã xóa
@Query("SELECT COALESCE(SUM(CASE WHEN t.type = ...TransactionType.income " +
       "THEN t.amount ELSE -t.amount END), 0) " +
       "FROM Transaction t WHERE t.account.id = :accountId AND t.deletedAt IS NULL")
BigDecimal computeBalanceForAccount(@Param("accountId") UUID accountId);
\end{lstlisting}

Endpoint \texttt{/api/accounts/\{id\}/reconcile} (và \texttt{/reconcile-all}) cho phép gọi tính
lại thủ công, đưa số dư về đúng tổng giao dịch -- kể cả khi số dư khởi tạo bị nhập lệch.

\subsection{Optimistic concurrency (\texttt{expectedUpdatedAt} $\rightarrow$ 409)}

Khi sửa giao dịch qua \texttt{PUT /api/transactions/\{id\}}, client có thể gửi
\texttt{expectedUpdatedAt}. Nếu bản trên server đã đổi so với mốc client kỳ vọng (thiết bị khác
vừa sửa), server ném \texttt{ConflictException} $\rightarrow$ \textbf{409}, \emph{không} ghi đè
âm thầm (chống lost-update):

\begin{lstlisting}[style=java,caption={TransactionService.updateTransaction -- optimistic concurrency check}]
// Optimistic concurrency: nếu bản trên server đã đổi so với expectedUpdatedAt -> 409
if (req.getExpectedUpdatedAt() != null && tx.getUpdatedAt() != null
        && !req.getExpectedUpdatedAt().isEqual(tx.getUpdatedAt())) {
    throw new ConflictException(
        "Giao dịch đã được thay đổi ở thiết bị khác. Hãy tải lại trước khi sửa.");
}
\end{lstlisting}

% ═════════════════════════════════════════════════════════════════════════════
\section{Cô lập dữ liệu đa tài khoản (4 lớp)}
\label{sec:isolation}

Một thiết bị có thể đăng nhập lần lượt nhiều tài khoản; dữ liệu giữa các tài khoản phải
tuyệt đối tách biệt. Hệ thống áp dụng bốn lớp phòng vệ:

\begin{table}[ht]
\centering
\caption{Bốn lớp cô lập dữ liệu giữa các tài khoản}
\label{tab:isolation}
\begin{tabular}{@{}lp{0.62\textwidth}@{}}
\toprule
\textbf{Lớp} & \textbf{Cơ chế} \\
\midrule
Lớp 1 -- Schema & Cột \texttt{user\_id} trên mọi entity; DAO lọc \texttt{WHERE user\_id = :id}. Backend luôn lấy \texttt{userId} từ \texttt{SecurityContext}, không tin tham số client. \\
Lớp 2 -- Write  & Stamp \texttt{userId} trước mọi insert (cả hành động người dùng lẫn ghi từ sync). \\
Lớp 3 -- Async  & ``Stale response guard'': bỏ qua response của user cũ nếu tài khoản đã đổi giữa lúc gửi và lúc nhận. \\
Lớp 4 -- Switch & Đổi tài khoản $\rightarrow$ \texttt{cancelAllRequests} + hủy worker + \texttt{clearAllTables} + \texttt{SyncPrefs.clearAll} $\rightarrow$ tạo \texttt{AppViewModel} mới. \\
\bottomrule
\end{tabular}
\end{table}

\noindent Phía Android, \texttt{TokenManager} (singleton) lưu token trong
\texttt{EncryptedSharedPreferences} (AES256-GCM, backed by Android Keystore) và cache trong RAM
để các luồng nền (OkHttp Dispatcher, WorkManager) không phải chạm Keystore vốn chậm. Tại tầng
mạng, \texttt{RetrofitClient} dùng chung một \texttt{Dispatcher}; khi đăng xuất gọi
\texttt{cancelAllRequests()} để hủy mọi request đang bay, tránh ``zombie response'' ghi dữ liệu
của user cũ vào DB của user mới.

% ═════════════════════════════════════════════════════════════════════════════
\section{Kiểm thử tích hợp end-to-end}
\label{sec:testing}

Lớp \texttt{MultiDeviceSessionFlowTest} boot toàn bộ ứng dụng trên \texttt{H2} in-memory
(profile \texttt{test}, \emph{không} đụng Neon production) và gọi API thật qua JDK
\texttt{HttpClient}. Một kịch bản kiểm chứng đồng thời cả ba mục: quản lý phiên đa thiết bị,
toàn vẹn số dư và optimistic concurrency. Các bước và kết quả mong đợi (đều PASS) tóm tắt ở
Bảng~\ref{tab:test}.

\begin{table}[ht]
\centering
\caption{Các bước trong MultiDeviceSessionFlowTest}
\label{tab:test}
\begin{tabular}{@{}clc@{}}
\toprule
\textbf{\#} & \textbf{Bước kiểm thử} & \textbf{Kết quả} \\
\midrule
1  & Đăng ký ``Thiết bị A'' (tự tạo phiên A)                       & 201 \\
2  & Đăng nhập ``Thiết bị B'' cùng tài khoản, song song           & 200 \\
3  & A liệt kê phiên $\rightarrow$ đúng 2, đúng 1 phiên ``current'' & 200 \\
4  & Xoay access token A bằng refresh token                       & 200 \\
5  & Dùng lại refresh token A cũ (đã xoay)                         & 401 \\
6  & A thu hồi phiên của B                                         & 204 \\
7  & Token B dùng ngay sau khi bị thu hồi (thu hồi tức thì)        & 401 \\
8  & Tạo ví số dư lệch 500k, không giao dịch $\rightarrow$ reconcile & 0 \\
9  & Thu 100k $\rightarrow$ số dư ví                               & 100k \\
10 & Chi 30k $\rightarrow$ số dư ví                                & 70k \\
11 & Xóa giao dịch chi $\rightarrow$ số dư ví                      & 100k \\
12 & Sửa với \texttt{expectedUpdatedAt} cũ (lost-update)          & 409 \\
13 & Sửa không kèm \texttt{expectedUpdatedAt} $\rightarrow$ 120k   & 120k \\
14 & Đăng xuất A $\rightarrow$ token A chết                        & 204/401 \\
\bottomrule
\end{tabular}
\end{table}

\noindent Chuỗi số dư \textbf{100k $\rightarrow$ 70k $\rightarrow$ 100k $\rightarrow$ 120k}
chứng minh nguyên tắc ``số dư $=\Sigma$ giao dịch'' hoạt động đúng qua create/expense/delete/update,
trong khi 409 ở bước 12 (kèm khẳng định số dư \emph{không} đổi) xác nhận cơ chế optimistic
concurrency chặn được lost-update từ thiết bị khác.

% ═════════════════════════════════════════════════════════════════════════════
\section{Tổng kết}
\label{sec:conclusion}

\begin{table}[ht]
\centering
\caption{Tóm tắt quyết định kỹ thuật}
\label{tab:summary}
\begin{tabular}{@{}lll@{}}
\toprule
\textbf{Quyết định} & \textbf{Lựa chọn} & \textbf{Lý do} \\
\midrule
Nguồn sự thật     & Room (UX) + PostgreSQL (hệ thống)   & tức thì + nhất quán toàn cục \\
Đồng bộ           & WorkManager + cursor delta (Slice)  & tiết kiệm băng thông, chịu lỗi \\
Xung đột          & Last-Write-Wins (\texttt{updatedAt})& đơn giản, hội tụ cho single-user \\
Toàn vẹn số dư    & $\Sigma$ giao dịch + 409             & chống trôi lệch, chống lost-update \\
Xác thực          & JWT stateless + \texttt{sid} + rotation & mở rộng tốt, thu hồi tức thì \\
Cô lập            & 4 lớp theo \texttt{user\_id}         & tách tuyệt đối dữ liệu đa tài khoản \\
\bottomrule
\end{tabular}
\end{table}

\noindent Các cơ chế trên phối hợp tạo nên một hệ thống offline-first vừa cho trải nghiệm tức
thì, vừa giữ nhất quán dữ liệu tài chính giữa nhiều thiết bị, với khả năng thu hồi phiên tức
thì và bảo đảm toàn vẹn số dư đã được kiểm chứng bằng kiểm thử tích hợp end-to-end.

\end{document}
